\begin{appendices}
\renewcommand{\theequation}{A\arabic{equation}}
\section{Model Codes}

\subsection{Code Model 1}
Deze vul ik even later in\dots

\subsection{Code Model 2}
    
Import Packages; generate grid; define constants:
\begin{lstlisting}
using Ferrite, SparseArrays, WriteVTK, Plots
name = "absorption_diffusion_convection_1d";

# Stop time, time step
dt = .005;
T = 1.3;
k = 2.0;         # rate of absorption
rho_sat = 1.0;   # saturation density of metal
vx = 1.5;        # x-velocity for convection
vy = 0.0;        # y-velocity for convection
velocity = vx;   # total velocity 
D = 0.1;         # Diffusivity for diffusion
n_grid = 100;    # number of grid points

# generates n_grid gridpoints between 0 and 1
grid = generate_grid(Line, (n_grid,), Vec((0.0,)), Vec((1.0,)))

# generate the same 
qr = QuadratureRule{RefLine}(2);

ip_rho = Lagrange{RefLine, 1}();
cellvalues_rho = CellValues(qr, ip_rho);

ip_rhos = Lagrange{RefLine, 1}()
cellvalues_rhos = CellValues(qr, ip_rhos);

# degree of freedom handler, 
dh = DofHandler(grid)
add!(dh, :rho, ip_rho)
add!(dh, :rho_s, ip_rhos)
close!(dh);
\end{lstlisting}
Constraint handler for boundary condition enforcement:
\begin{lstlisting}
# for now we have no constraints
ch = ConstraintHandler(dh)

#create left_slit
dbc_left = Dirichlet(
    :rho,
    getfacetset(grid, "left"),
    (x, t) -> 1.0
)

add!(ch, dbc_left)
close!(ch)
update!(ch, 0.0)
\end{lstlisting}
\begin{equation}
    M_{ij} = \int_{\Omega} \phi_i \phi_j d\Omega
\end{equation}
\begin{lstlisting}
function doassemble_M!(M, cellvalues_rho, cellvalues_rhos, dh)
    assembler = start_assemble(M)

    # number of basis functions per field
    n_basefuncs_rho  = getnbasefunctions(cellvalues_rho)
    n_basefuncs_rhos = getnbasefunctions(cellvalues_rhos)

    # loop over all cells
    for cell in CellIterator(dh)
        # get global DOFs for this cell
        cell_dofs = celldofs(cell)
        ndofs_cell = length(cell_dofs)
        Me = zeros(ndofs_cell, ndofs_cell)

        # rho
        reinit!(cellvalues_rho, cell)
        for q in 1:getnquadpoints(cellvalues_rho)
            dx = getdetJdV(cellvalues_rho, q)
            for i in 1:n_basefuncs_rho
                psi_i = shape_value(cellvalues_rho, q, i)
                for j in 1:n_basefuncs_rho
                    psi_j = shape_value(cellvalues_rho, q, j)
                    # :rho DOFs are first in cell_dofs
                    Me[i, j] += psi_i * psi_j * dx
                end
            end
        end

        # rhos
        reinit!(cellvalues_rhos, cell)
        for q in 1:getnquadpoints(cellvalues_rhos)
            dx = getdetJdV(cellvalues_rhos, q)
            for i in 1:n_basefuncs_rhos
                psi_i = shape_value(cellvalues_rhos, q, i)
                for j in 1:n_basefuncs_rhos
                    psi_j = shape_value(cellvalues_rhos, q, j)
                    # :rho_s DOFs are after :rho in cell_dofs
                    offset = n_basefuncs_rho
                    Me[offset + i, offset + j] += psi_i * psi_j * dx
                end
            end
        end

        assemble!(assembler, cell_dofs, Me)
    end

    return M
end
\end{lstlisting}
\begin{equation}
    K_{ij} = \int_{\Omega} \nabla \phi_i \cdot \nabla \phi_j d\Omega
\end{equation}
\begin{lstlisting}
function doassemble_K!(
    K::SparseMatrixCSC, 
    cellvalues_rho::CellValues, 
    cellvalues_rhos::CellValues, 
    dh::DofHandler
    )

    assembler = start_assemble(K)

    # number of basis functions per field
    n_basefuncs_rho  = getnbasefunctions(cellvalues_rho)
    n_basefuncs_rhos = getnbasefunctions(cellvalues_rhos)

    for cell in CellIterator(dh)

        # get global DOFs for this cell
        cell_dofs = celldofs(cell)
        ndofs_cell = length(cell_dofs)
        Ke = zeros(ndofs_cell, ndofs_cell)

        # rho is the only important bit since rho_s does not diffuse
        reinit!(cellvalues_rho, cell)

        for q_point in 1:getnquadpoints(cellvalues_rho)
            dx = getdetJdV(cellvalues_rho, q_point)

            for i in 1:n_basefuncs_rho
                dpsi_i = shape_gradient(cellvalues_rho, q_point, i)

                for j in 1:n_basefuncs_rho
                    dpsi_j = shape_gradient(cellvalues_rho, q_point, j)
                    # D∫ (∇ψ^i) ⋅ (∇ψ^j)dx
                    Ke[i, j] += D * (dpsi_j ⋅ dpsi_i) * dx
                end
            end
        end

        assemble!(assembler, cell_dofs, Ke)
    end

    return K
end
\end{lstlisting}
\begin{equation}
    C_{ij} = -v\int_{\Omega} \phi_i \nabla \phi_j d\Omega
\end{equation}
\begin{lstlisting}
function doassemble_C!(
    K::SparseMatrixCSC, 
    cellvalues_rho::CellValues, 
    cellvalues_rhos::CellValues, 
    dh::DofHandler
    )

    n_basefuncs_rho = getnbasefunctions(cellvalues_rho)
    Ce = zeros(n_basefuncs_rho, n_basefuncs_rho)

    assembler = start_assemble(C)

    for cell in CellIterator(dh)
        
        # get global DOFs for this cell
        cell_dofs = celldofs(cell)
        ndofs_cell = length(cell_dofs)
        Ce = zeros(ndofs_cell, ndofs_cell)

        # rho is the only important bit since rho_s does not diffuse
        reinit!(cellvalues_rho, cell)
        fill!(Ce, 0)

        for q_point in 1:getnquadpoints(cellvalues_rho)
            dx = getdetJdV(cellvalues_rho, q_point)

            for i in 1:n_basefuncs_rho
                psi_i = shape_value(cellvalues_rho, q_point, i)

                for j in 1:n_basefuncs_rho
                    dpsi_j = shape_gradient(cellvalues_rho, q_point, j)
                    # -v \psi^i(\partial_x\psi^j)dx
                    Ce[i, j] += (velocity ⋅ dpsi_j) *  psi_i * dx
                end
            end
        end
        assemble!(assembler, cell_dofs, Ce)
    end

    return C
end    
\end{lstlisting}
Generate the mass, stiffness and convection matrices:
\begin{lstlisting}
M = allocate_matrix(dh);
K = allocate_matrix(dh);
C = allocate_matrix(dh);

M = doassemble_M!(M, cellvalues_rho, cellvalues_rhos, dh);
K = doassemble_K!(K, cellvalues_rho, cellvalues_rhos, dh);
C = doassemble_C!(C, cellvalues_rho, cellvalues_rhos, dh);
\end{lstlisting}
For the ImEx method:
\begin{equation}
    \underbrace{\left(M + \Delta t(K+C)\right)}_{=A} \rho^{n+1} = M \rho^n + M\dot{\mathbf{m}}_h^n(\mathbf{u}^n)\Delta t
\end{equation}
Also we define the LU factorization of the left-hand side matrix for faster computation of the inverse:
\begin{lstlisting}
A = M + dt .* (K+C);
A_fact = lu(A);
\end{lstlisting}
Make the right-hand side vector; apply bcs; define global dofs for $\rho$ and $\rho_s$ for mass absorption calculation:
\begin{lstlisting}
ndofs_total = ndofs(dh)
u_n = zeros(ndofs_total)
apply!(u_n, ch)

rhsdata = get_rhs_data(ch, A);
apply!(A, ch);

# find global dofs for rho and rho_s for mass absorption calculation
rho_dofs  = Int[]
rhos_dofs = Int[]

for cell in CellIterator(dh)
    cdofs = celldofs(cell)

    append!(rho_dofs,  cdofs[dof_range(dh, :rho)])
    append!(rhos_dofs, cdofs[dof_range(dh, :rho_s)])
end

rho_dofs  = unique(rho_dofs);
rhos_dofs = unique(rhos_dofs);

# for savind and plotting
steps_n = length(dt:dt:T)
u_hist = zeros(steps_n, ndofs_total)
\end{lstlisting}
Main time loop:
\begin{lstlisting}
mdot = zeros(ndofs_total)
# A_fact = lu(A)

# finally the timeloop
for (step, t) in enumerate(dt:dt:T)
    
    #First of all, we need to update the Dirichlet boundary condition values.
    update!(ch, t)


    # now we reset ̇m.
    mdot = zeros(ndofs_total)

    # important: because of our assumptions above, we iterate over degrees of freedom/nodes, not cells
    # the cell relavance will be added after M*mdot
    for (i_rho, i_rhos) in zip(rho_dofs, rhos_dofs)

        # for now: mdto = kρ(ρsat-ρs)
        md = k * u_n[i_rho] * (rho_sat - u_n[i_rhos])

        mdot[i_rhos] += md
        mdot[i_rho]  -= md
    end

    # b without boundary elements
    b = M * u_n  + dt * (M * mdot)

    #Then, we can apply the boundary conditions of the current time step.
    apply_rhs!(rhsdata, b, ch)

    #Finally, we can solve the time step and save the solution afterwards.
    u = A \ b

    
    VTKGridFile(joinpath(name, "$(name)-$step"), dh) do vtk
        write_solution(vtk, dh, u)
        pvd[t] = vtk
    end
    #At the end of the time loop, we set the previous solution to the current one and go to the next time step.
    u_n .= u
    u_hist[step, :] .= u_n

end    
\end{lstlisting}

\subsection{Code Model 3}

\subsection{Code Model 4}

\subsection{Code Model 5}

\end{appendices}